\subsection{Functional calculus in unital stellar algebras}

In this number, A is a unital stellar algebra and x is a normal element of A.

Let B be the unital stellar subalgebra of A generated by x; it is commutative and contained in the bicommutant of \(\{x, x^{*}\}\) , hence in the bicommutant of x (corollary of Proposition 6 in I, p.~106). The Gelfand transform \(\mathcal{G}_{\mathrm{B}}: \mathrm{B} \to \mathcal{C}(\mathrm{X}(\mathrm{B}))\) is an isomorphism of stellar algebras (th. 1 of I, p.~108). We have \(\mathrm{Sp}_{\mathrm{B}}(x) = \mathrm{Sp}_{\mathrm{A}}(x)\) (prop. 5 of I, p.~106).

Lemma 10. --- The map \(\mathrm{ev}_x\colon \chi \mapsto \chi (x)\) induces a homeomorphism of \(\mathsf{X}(\mathbf{B})\) onto \(\mathrm{Sp}_{\mathrm{A}}(x)\) .

The map \(x \mapsto \chi(x)\) of \(\mathsf{X}(\mathsf{B})\) into \(\mathbf{C}\) is continuous, and its image is equal to \(\mathrm{Sp}_{\mathsf{B}}(x)\) by prop. 6 of I, p.~37, hence to \(\mathrm{Sp}_{\mathsf{A}}(x)\) . Since the characters of B are Hermitian (lemma 9 of I, p.~107), characters of B that agree on \(x\) are also equal on \(x^{*}\) and therefore are equal on the unital stellar subalgebra B generated by \(x\) . This proves that the map \(\mathrm{ev}_x\) is injective. Since \(\mathsf{X}(\mathsf{B})\) is compact and \(\mathbf{C}\) is separated, it induces a homeomorphism from \(\mathsf{X}(\mathsf{B})\) onto its image, hence the lemma.

From the lemma we deduce an isomorphism of stellar algebras \(\varphi_x\colon \mathcal{C}(\mathrm{Sp}_{\mathrm{A}}(x))\to \mathcal{C}(\mathrm{X(B)})\) . It maps a function \(f\in \mathcal{C}(\mathrm{Sp}_{\mathrm{A}}(x))\) to the function \(\chi \mapsto f(\chi (x))\) . The map \(\mathcal{G}_{\mathrm{B}}^{-1}\circ \varphi_{x}\) is an isomorphism of the stellar algebra \(\mathcal{C}(\mathrm{Sp}_{\mathrm{A}}(x))\) onto B.

DEFINITION 5. --- The involutive morphism \(f \mapsto (\mathcal{G}_{\mathrm{B}}^{-1} \circ \varphi_x)(f)\) of \(\mathcal{C}(\mathrm{Sp}_{\mathrm{A}}(x))\) to A is called the continuous functional calculus map of \(x\) in A. We denote it \(f \mapsto f(x)\) .

Remark. --- The map \(f \mapsto f(x)\) is isometric; its image is the unital stellar subalgebra B generated by x, which is contained in the bicommutant of x.

If \(f\) is the restriction to \(\mathrm{Sp}_{\mathrm{A}}(x)\) of a function of the form \(z \mapsto \mathrm{P}(z, \overline{z})\) , where \(\mathrm{P} \in \mathbf{C}[\mathrm{X}, \mathrm{Y}]\) is a polynomial, then we have \(f(x) = \mathrm{P}(x, x^{*})\) in the usual algebraic sense.

Examples. --- 1) Suppose that there exists \(\lambda \in \mathbf{C}\) such that \(\mathrm{Sp}_{\mathrm{A}}(x)\) is reduced to \(\lambda\) . We then have \(x = \lambda \cdot 1\) . Indeed, the identity function of \(\mathrm{Sp}_{\mathrm{A}}(x)\) is equal to \(\lambda\) , hence its image under the functional calculus map, that is, \(x\) is equal to \(\lambda \cdot 1\) .

\begin{enumerate}
\def\labelenumi{\arabic{enumi})}
\setcounter{enumi}{1}
\item
  For \(x\) to be Hermitian, it is necessary and sufficient that \(\mathrm{Sp}_{\mathrm{A}}(x)\) be contained in \(\mathbf{R}\) . Indeed, let \(f\) be the continuous map on \(\mathrm{Sp}_{\mathrm{A}}(x)\) given by \(f(z) = z - \overline{z}\) . Then \(x\) is Hermitian if and only if \(f(x) = 0\) , that is, if \(f\) is zero, that is, if \(\mathrm{Sp}_{\mathrm{A}}(x)\) is contained in \(\mathbf{R}\) .
\item
  For \(x\) to be unitary, it is necessary and sufficient that its spectrum be contained in the unit circle of \(\mathbf{C}\) . Indeed, let \(f \in \mathcal{C}(\mathrm{Sp}_{\mathrm{A}}(x))\) be the function defined by \(f(z) = z\overline{z} - 1\) ; the element \(x\) is unitary if and only if \(f(x) = 0\) , that is, if \(f\) is zero.
\end{enumerate}

Proposition 7. — The map \(f \mapsto f(x)\) is the unique unital morphism of involutive algebras from \(\mathcal{C}(\mathrm{Sp}_{\mathrm{A}}(x))\) into A such that the identity map \(z\) of \(\mathrm{Sp}_{\mathrm{A}}(x)\) has image \(x\) .

Indeed, the unital subalgebra of \(\mathcal{C}(\mathrm{Sp}_{\mathrm{A}}(x))\) generated by the elements \(z\) and \(\overline{z}\) of \(\mathcal{C}(\mathrm{Sp}_{\mathrm{A}}(x))\) is dense in \(\mathcal{C}(\mathrm{Sp}_{\mathrm{A}}(x))\) (TG, X, p.~40, cor. 1). Since every morphism of involutive algebras from \(\mathcal{C}(\mathrm{Sp}_{\mathrm{A}}(x))\) to A is continuous (I, p.~104, prop. 2), there exists at most one morphism of involutive algebras from \(\mathcal{C}(\mathrm{Sp}_{\mathrm{A}}(x))\) to A which maps \(z\) to \(x\) .

The following corollary shows that when \(f\) is the restriction of a holomorphic function in a neighborhood of \(\mathrm{Sp}_{\mathrm{A}}(x)\) the definition of \(f(x)\) coincides with that of the holomorphic functional calculus in one variable from no. 9 of I, p.~74.

COROLLARY 1. --- Let \(f \in \mathcal{O}(\mathrm{Sp}_{\mathrm{A}}(x))\) be a germ of a holomorphic function in a neighborhood of \(\mathrm{Sp}_{\mathrm{A}}(x)\) and let \(\widetilde{f} \in \mathcal{C}(\mathrm{Sp}_{\mathrm{A}}(x))\) be the continuous function on \(\mathrm{Sp}_{\mathrm{A}}(x)\) associated with \(f\) . We have \(\widetilde{f}(x) = f(x)\) , where \(f(x)\) is the element of A given by the holomorphic functional calculus.

Indeed, the map \(f \mapsto \widetilde{f}(x)\) is a unital continuous morphism from \(\mathcal{O}(\mathrm{Sp}_{\mathrm{A}}(x))\) to A which maps the germ of the identity function in a neighborhood of \(\mathrm{Sp}_{\mathrm{A}}(x)\) to \(x\) . The result then follows from th. 5 of I, p.~74.

COROLLARY 2. --- Let \(f \in \mathcal{C}(\mathrm{Sp}_{\mathrm{A}}(x))\) .

\begin{enumerate}
\def\labelenumi{\alph{enumi})}
\tightlist
\item
  We have
\end{enumerate}

\[
\operatorname{Sp} _ {\mathrm{A}} (f (x)) = f (\operatorname{Sp} _ {\mathrm{A}} (x));
\]

\begin{enumerate}
\def\labelenumi{\alph{enumi})}
\setcounter{enumi}{1}
\tightlist
\item
  For every \(g \in \mathcal{C}(\mathrm{Sp}_{\mathrm{A}}(f(x)))\) , one has \((g \circ f)(x) = g(f(x))\) .
\end{enumerate}

Since \(f(x)\) belongs to the full subalgebra B of A, we have \(\mathrm{Sp}_{\mathrm{A}}(f(x)) = \mathrm{Sp}_{\mathrm{B}}(f(x))\) (prop. 5 of I, p.~106). The isomorphism \(f \mapsto f(x)\) of \(\mathcal{C}(\mathrm{Sp}_{\mathrm{A}}(x))\) to B preserves the spectrum; we therefore have \(\mathrm{Sp}_{\mathrm{B}}(f(x)) = \mathrm{Sp}_{\mathcal{C}(\mathrm{Sp}_{\mathrm{A}}(x))}(f) = f(\mathrm{Sp}_{\mathrm{A}}(x))\) (example 3 of I, p.~17). This proves assertion \(a\) ).

The map \(g \mapsto (g \circ f)(x)\) is a unital morphism of involutive algebras from \(\mathcal{C}(\mathrm{Sp}_{\mathrm{A}}(f(x)))\) to A which maps the identity function of \(\mathrm{Sp}_{\mathrm{A}}(f(x))\) to \(f(x)\) . By prop. 7, we therefore have \((g \circ f)(x) = g(f(x))\) for all \(g \in \mathcal{C}(\mathrm{Sp}_{\mathrm{A}}(f(x)))\) .

Example 4. --- Let X be a locally compact space and let A = \(\mathcal{C}_{b}(X)\) be the commutative unital stellar algebra of continuous bounded functions on X (example 2 of I, p.~102). Let \(g \in A\) ; its spectrum S is the closure in C of the set \(g(\mathrm{X})\) of the values of g (example 3 of I, p.~17). The functional calculus mapping of g is then the mapping \(f \mapsto f \circ g\) , for \(f \in \mathcal{C}(\mathrm{S})\) . Indeed, this map is a unital morphism of stellar algebras such that the identity map of S has image g.

In the case where X is compact, we have \(A = \mathcal{C}(X)\) and \(S = g(X)\) .

Let \(\pi\colon\mathrm{A}\to\mathrm{A}^{\prime}\) be a unital morphism of unital stellar algebras. The element \(\pi(x)\) of \(\mathrm{A}^{\prime}\) is normal and its spectrum relative to \(\mathrm{A}^{\prime}\) is contained in \(\mathrm{Sp}_{\mathrm{A}}(x)\) . We then have:

PROPOSITION 8. --- Let \(f \in \mathcal{C}(\mathrm{Sp}_{\mathrm{A}}(x))\) . Still denoting \(f\) the restriction of \(f\) to \(\mathrm{Sp}_{\mathrm{A}'}(\pi(x))\) , we have the equality \(\pi(f(x)) = f(\pi(x))\) . In particular, for every \(\chi \in \mathsf{X}(\mathrm{A})\) , one has \(\chi(f(x)) = f(\chi(x))\) .

Let \(z\) be the identity map of \(\mathrm{Sp}_{\mathrm{A}}(x)\) . The maps defined by \(f \mapsto \pi(f(x))\) and \(f \mapsto f(\pi(x))\) are continuous unital morphisms of involutive algebras from \(\mathcal{C}(\mathrm{Sp}_{\mathrm{A}}(x))\) to B that map \(z\) to \(\pi(x)\) . These morphisms therefore coincide on the unital involutive subalgebra of \(\mathcal{C}(\mathrm{Sp}_{\mathrm{A}}(x))\) generated by \(z\) . Since this one is dense in \(\mathcal{C}(\mathrm{Sp}_{\mathrm{A}}(x))\) (TG, X, p.~40, cor. 1), these morphisms are equal.

COROLLARY. --- Suppose that A is commutative. For every \(f \in \mathcal{C}(\mathrm{Sp}_{\mathrm{A}}(x))\) , one has \(\mathcal{G}_{\mathrm{A}}(f(x)) = f \circ \mathcal{G}_{\mathrm{A}}(x)\) .

It suffices to apply Prop. 8 to the Gelfand transform \(\mathcal{G}_{\mathrm{A}}\) from A to \(\mathcal{C}(\mathsf{X}(\mathrm{A}))\) and to note (example above) that \(f(\mathcal{G}_{\mathrm{A}}(x)) = f \circ \mathcal{G}_{\mathrm{A}}(x)\) .

